\chapter{语言大模型}
\begin{quote}\small
\textit{本章摘要。}语言大模型首先是条件概率模型：它根据已有文本预测下一个词元。要把这种预测能力用于设备维护，还须解决三个不同的问题：如何根据示范和偏好调整回答，如何从现行手册取得依据，以及如何检验生成结果。本章从词元表示与因果建模出发，依次讨论训练、解码和外部证据的使用。
\end{quote}
图像模型可以指出裂纹的位置，图模型可以利用设备之间的连接关系。然而，值班人员据此安排复检时，还得查明该型号允许的偏差、适用的工况，以及手册的现行版本。这些条件通常写在文本中，语言模型因而有了明确的用途：从文档和记录中组织与当前问题相关的说明。

这里必须分清能力的来源。预训练赋予模型续写文本的能力；微调与偏好优化使它更常按任务要求作答；检索在推理时补充具体条款。前两者改变模型参数，检索通常只改变本次输入。以往工单可以帮助训练，却不能作为尚未发生的维修事件的现场证据。本章将始终用这一时间边界检验模型的回答。
\section{词元、表示与语言建模}

语言模型处理的是词元（token）序列，而不是直接处理一串汉字。分词器（tokenizer）按既定词表把文本切成词元，再赋予编号；词元可能是一个字、词的一部分、标点或字节片段。为便于计算，假设“温度升高”被切成“温度”和“升高”，编号分别为$3$和$7$。编号只标识词表中的位置，不表示语义大小。模型通过嵌入矩阵把编号转换为可训练的向量；实际切分结果则由所用分词器决定。

先看数据怎样进入网络。一个批次若有两段文本，每段补齐到四个词元，输入编号的形状就是$2\times4$。嵌入宽度若为八，查表后的形状是$2\times4\times8$。设词表共有十项，网络会为每个位置给出十个候选词元的得分，形成$2\times4\times10$的输出。得分称为 logit；沿词表维度施加 softmax 后，才得到该位置上和为一的概率分布。

自回归模型只依据已知前缀预测后续词元。一段序列的概率因此由条件概率连乘得到：
\begin{equation}
 p_\theta(x_{1:T})=\prod_{t=1}^{T}p_\theta(x_t\mid x_{<t}),\label{eq:llm-chain-probability}
\end{equation}
其中$x_{<t}$表示第$t$个词元之前的全部内容。这个等式来自概率的链式法则，并没有假设词元相互独立。训练时，模型要使语料中实际出现的词元获得较高概率；生成时，它则从当前分布中选取词元，接到前缀后继续预测。

从编号到概率还需一个计算网络。本章采用 Transformer：词嵌入与位置信息构成初始表示，多层因果自注意力和前馈网络更新表示，最后映射到词表得分\cite{vaswani2017attention}。先看训练样本的输入与目标，再展开这些计算。

\paragraph{一段文本怎样产生训练标签。}
仍以“温度升高”为例，在序列两端加入起始符与结束符。训练输入为“起始符、温度、升高”，对应目标为“温度、升高、结束符”。第一个位置只凭起始符预测“温度”；第二个位置再利用“温度”预测“升高”；第三个位置预测何时结束。若模型给三个目标的概率依次为$1/2,1/4,1/2$，整段序列的概率就是$1/16$，平均负对数似然为$\log16/3$，困惑度为$16^{1/3}\approx2.52$。这里及下文的对数均为自然对数。

训练时，上述三个目标可以在一次前向计算中并行计分，因为真实前缀已经给定，因果掩码又阻止每个位置偷看后文。这种使用真实前缀的做法称为教师强制。生成时并不知道下一个词元，只能选出一个、接到前缀末尾，再计算下一步分布。两种过程使用同一模型，但生成中的早期错误会进入后续预测的条件。

\paragraph{链式分解如何变成可计算的训练目标。}
现在写出单个位置的损失。设词表大小为$V_{\rm voc}$，$x_t\in\{1,\ldots,V_{\rm voc}\}$，嵌入矩阵$E\in\R^{V_{\rm voc}\times d}$。起始符（BOS）为第一个词元提供条件，结束符（EOS）使模型学会停止。令因果网络从前缀得到$h_t\in\R^d$，输出矩阵$U\in\R^{V_{\rm voc}\times d}$。预测分布、负对数似然及其对得分的导数依次为
\begin{gather}
 z_t=Uh_t+b,\qquad p_{t,j}=\frac{e^{z_{t,j}}}{\sum_{k=1}^{V_{\rm voc}}e^{z_{t,k}}},\\
 \ell_t=-\log p_{t,x_t}=-z_{t,x_t}+\log\sum_ke^{z_{t,k}},\\
 \frac{\partial\ell_t}{\partial z_{t,j}}=p_{t,j}-\mathbf1_{j=x_t}.
\end{gather}
最后一个等式说明了交叉熵的更新方向：对真实词元，梯度为$p_{t,x_t}-1$；对其余词元，梯度为$p_{t,j}$。若得分可以独立调整，梯度下降便会提高真实词元的相对得分。实际更新作用于共享网络参数，各位置的梯度会相互叠加，故单步更新不保证每个样本的概率都改善。梯度传向隐藏状态时，有
\begin{equation}
 \frac{\partial\ell_t}{\partial h_t}=U^{\mathsf T}(p_t-e_{x_t}),
\end{equation}
其中$p_t\in\R^{V_{\rm voc}}$为预测概率向量，$e_{x_t}$在真实词元处为一、其他位置为零；$\mathbf1_{j=x_t}$是对应的指示量。由此得到的$d$维梯度继续向前传播。计算对数求和指数时，常取$m=\max_k z_k$，写成$\log\sum_ke^{z_k}=m+\log\sum_ke^{z_k-m}$，以免直接计算指数时溢出。

实现时须区分两种掩码。因果注意力掩码规定某个位置\emph{能读取}哪些词元，损失掩码规定哪些输出\emph{参与计分}。例如指令微调可以只在回答位置计损失，同时让回答读取整段问题。填充位置不计分，损失应按实际计分的词元数归一化。此外，部分训练接口会自动把标签右移；若又手工移位，输入和目标就会错位。无论采用何种接口，目标$x_t$都不能出现在预测它的可见前缀中。

固定输入前缀$c$，设训练数据中下一个词元的真实条件分布为$q(\cdot\mid c)$，则期望交叉熵分解为
\begin{equation}
 -\sum_jq_j\log p_j=-\sum_jq_j\log q_j+\KL(q\|p).
\end{equation}
右侧第一项是数据分布自身的熵，与模型参数无关；训练所能降低的是第二项，即$q$与$p$之间的差异。这解释了最大似然训练的能力，也指出了它的边界：模型学会的是语料中的用语和事实分布，而不是“哪些规程今天仍适用”的判定程序。若语料保留了旧版条款，模型把旧条款续写得更加熟练也符合训练目标。现行依据必须另行查证。

词表的设计还会影响计算成本。词表较小，罕见字符串往往被拆成更多词元，序列随之变长；词表较大，嵌入与输出层的参数增多，低频词元也更难充分训练。设备编号、故障码和计量单位尤其容易被切碎，部署时应检验模型能否完整抄录。由于标准自注意力的计算量随序列长度近似二次增长，估计上下文成本应数词元，而不是数字符。

因果 Transformer 怎样利用前缀？设某层输入为矩阵$H$。注意力先把它投影为查询、键和值，再由查询与键的相似度决定读取各个值的权重：
\begin{equation}
 Q=HW_Q,\quad K=HW_K,\quad V=HW_V,
 \qquad A=\operatorname{softmax}\left(\frac{QK^{\mathsf T}}{\sqrt{d_k}}+M\right),
\end{equation}
其中$d_k$是键的宽度，softmax 按行计算。掩码$M$在$j>i$时取$-\infty$，其他位置取零，于是第$i$个位置只会给自己和此前位置分配权重。这是“因果”的含义：训练时可以并行算出所有位置的输出，但每个输出仍服从从左到右的信息约束。

若序列长$T$、表示宽$d$，则$H\in\R^{T\times d}$，$W_Q,W_K\in\R^{d\times d_k}$，$W_V\in\R^{d\times d_v}$。因此权重矩阵$A\in\R^{T\times T}$，读出的内容$AV\in\R^{T\times d_v}$。例如某一行仅有两个可见位置，得分为$0$与$\log3$，对应权重便是$1/4$与$3/4$：模型对两个位置加权读取，而非只挑一个。多头注意力使用不同投影，拼接后还需映射回宽度$d$；前馈层逐位置加工表示，残差与归一化帮助深层网络训练。注意力头可能形成不同分工，但不能仅凭头的编号赋予它固定语义。

归一化与残差的次序也是结构选择。预归一化把归一化放在注意力或前馈层之前，残差分支因此保留较直接的梯度通路；后归一化则在相加之后处理表示。图\ref{fig:llm-computation-chain}概括了这些层怎样把词元序列变成词表得分，箭头表示前向计算，图中的因果注意力仍须遵守上述可见性约束。

\begin{figure}[H]
\centering
\begin{tikzpicture}[>=Latex, node distance=0.65cm, every node/.style={font=\small}]
 \node[draw, rounded corners, fill=blue!10, minimum width=3cm, minimum height=0.7cm] (tok) {Token 序列};
 \node[draw, rounded corners, fill=orange!12, below=of tok, minimum width=3cm, minimum height=0.7cm] (emb) {词嵌入 + 位置};
 \node[draw, rounded corners, fill=green!12, below=of emb, minimum width=3cm, minimum height=0.7cm] (att) {多头因果注意力};
 \node[draw, rounded corners, fill=yellow!20, below=of att, minimum width=3cm, minimum height=0.7cm] (ffn) {前馈网络 + 残差};
 \node[draw, rounded corners, fill=red!10, below=of ffn, minimum width=3cm, minimum height=0.7cm] (head) {词表 logits};
 \draw[->, thick] (tok)--(emb); \draw[->, thick] (emb)--(att); \draw[->, thick] (att)--(ffn); \draw[->, thick] (ffn)--(head);
\end{tikzpicture}
\caption{自回归语言模型从词元序列到词表得分的计算。箭头表示前向计算的顺序；因果注意力使每个位置只能使用可见前缀。}
\label{fig:llm-computation-chain}
\end{figure}

\section{预训练和后训练}
\subsection{架构选择与预训练数据}
前面分析的是因果语言模型，但 Transformer 并非只有一种信息可见方式。根据任务是要理解一段已有文本，还是要逐词生成新文本，模型通常采用不同结构\cite{devlin2019bert,brown2020gpt3}。BERT 类编码器使用双向注意力，通过恢复被遮蔽的词元学习上下文表示，适合分类与检索重排序；GPT、LLaMA 类解码器采用因果注意力，直接训练下一词元预测；T5 类模型先用编码器读取完整输入，再由解码器条件生成输出，适合摘要和文本转换。

以“温度［遮蔽］，请复检”为例，BERT 可以同时利用遮蔽处左右两侧的文字；因果解码器只能依据左侧前缀；T5 则把整个输入交给编码器，再生成所需输出。这种差异是模型内部的信息约束，不等同于现实任务的时间约束。即便使用双向编码器，在线预警也只能输入预测时刻已经取得的记录。

decoder-only 模型对整段序列优化
\begin{equation}
 \mathcal{L}_{\mathrm{CLM}}=-\sum_{t=1}^{T}\log p_\theta(x_t\mid x_{<t}),
\end{equation}
BERT 则只在被遮蔽位置计算预测损失。两者见到的上下文和接受的监督不同，不能仅用参数量比较其任务适用性。例如处理设备手册时，编码器可以负责检索条款，生成模型再据条款写出说明。这样可分别衡量“找到了什么”和“回答是否忠于所找内容”。

混合专家（Mixture-of-Experts，MoE）处理的是另一维度的问题：怎样增加模型容量，而不让每个词元都经过全部参数。稀疏 MoE 用路由器为当前词元表示选择少数专家网络：
\begin{equation}
 y=\sum_{e\in\mathcal{S}(x)}g_e(x)E_e(x),\qquad |\mathcal{S}(x)|\ll N_{\rm exp}.
\end{equation}
其中$N_{\rm exp}$为专家总数，$\mathcal S(x)$为当前选中的索引集合，$E_e(x)$与$g_e(x)$分别是专家输出和路由权重。各输出须同维，才能加权求和。这里的“专家”只是可训练子网络；MoE 往往替换 Transformer 中的部分前馈层，可与不同注意力结构结合。

例如 top-1 路由每次选择一个专家，top-2 路由选择两个。总参数量虽可很大，单次计算只激活其中一部分；但路由不均会造成某些专家拥堵，跨设备分配还会产生通信开销。因此报告 MoE 成本时，应区分模型存储、每词元激活计算与设备间通信。

\begin{table}[htbp]
\centering
\caption{语言模型架构与工业任务匹配}
\label{tab:llm-architectures}
\scriptsize
\setlength{\tabcolsep}{3pt}
\resizebox{\linewidth}{!}{%
\begin{tabular}{p{2.7cm}p{3.8cm}p{3.7cm}p{3.8cm}}
\hline
架构路线 & 训练与注意力方向 & 优势 & 工业任务和风险 \\
\hline
BERT 类 encoder-only & 双向掩码建模 & 表示稳定、分类和检索高效 & 事件分类、实体抽取、重排序；不直接生成长文本 \\
GPT/LLaMA 类 decoder-only & 因果下一 token 预测 & 生成、对话和工具调用自然 & 报告草拟、脚本生成；需防幻觉和提示注入 \\
T5 类 encoder--decoder & 条件文本到文本 & 输入输出长度灵活、任务统一 & 摘要、规范化和文本转换；部署链路更复杂 \\
MoE 语言模型 & 路由器选择少数专家 & 总容量大而激活成本可控 & 多领域共享；需监控负载和通信 \\
\hline
\end{tabular}}
\end{table}

表\ref{tab:llm-architectures}的前三行按信息可见性与输出形式比较架构，最后一行则给出可叠加的容量设计。选型时应先确定要做分类、检索还是生成，再评估模型大小与部署成本。

对于因果语言模型，预训练最直接的指标是平均负对数似然：
\begin{equation}
 \mathcal{L}_{LM}=-\frac1T\sum_{t=1}^{T}\log p_\theta(x_t\mid x_{<t}).
\end{equation}
将平均负对数似然取指数，便得到困惑度$\operatorname{PPL}=\exp(\mathcal{L}_{LM})$。若每一步都在$k$个等可能词元中猜测，困惑度恰为$k$；一般分布下，它反映平均预测难度。分词器、测试集或计分位置一旦改变，困惑度就不宜直接比较；即使可比，它仍只衡量文本预测，不衡量规程是否正确。

训练效果受参数量、有效训练词元数和计算预算共同影响。参数增加而数据仍高度重复，模型可能更擅长记忆模板，却未见到更多工况。因而预训练语料需要去重、处理敏感信息，并与评测集隔离；工业文档还须保存设备型号、版本和生效时间。

比较训练实验时，除参数量外，还应记录有效词元数、序列长度、训练步数及去重比例。若评测问题或答案模板混入训练集，测试分数将高估泛化能力。旧版规程不一定毫无价值，但必须连同来源、许可和曾经适用的时间范围一起保存，以免在当前问答中被误用。

\subsection{后训练}
预训练让模型掌握词元之间的统计关系，但并未规定它应怎样回应维修人员的提问。后训练可利用两类监督：示范回答告诉模型“怎样写”，成对偏好告诉模型“两种写法中哪一种更可取”。InstructGPT 使用奖励模型和策略优化处理人类反馈\cite{ouyang2022instruct}；直接偏好优化（DPO）则从回答对构造语言模型的训练目标\cite{rafailov2023dpo}。无论采用哪条路线，事实正确性仍取决于数据和独立检验。

监督微调（SFT）以问题和示范回答为样本，继续最小化回答词元的负对数似然。它可以更新全部参数，也可以只更新一小部分。低秩适配（LoRA）属于后一种做法：冻结原权重$W$，只训练一个低秩增量\cite{hu2021lora}：
\begin{equation}
 W'=W+\frac{\alpha}{r}BA,\qquad B\in\mathbb{R}^{d\times r},\ A\in\mathbb{R}^{r\times k},\ r\ll\min(d,k).
\end{equation}
这里$W\in\R^{d\times k}$，$r$控制增量的秩，$\alpha/r$控制其尺度。训练完成后，只需保存小矩阵$A,B$，即可在同一基座上管理不同任务的适配器；但每个适配器仍须注明所配套的基座版本和适用任务。

对一个$d\times k$矩阵，全量更新需要$dk$个参数，LoRA 仅需$r(d+k)$个，参数节省比例为
\begin{equation}
 1-\frac{r(d+k)}{dk}.
\end{equation}
节省参数的代价是限制更新方向。若任务所需变化难以由秩$r$的矩阵表达，性能可能低于全量微调。评价 LoRA 时，任务效果、训练显存和推理开销须与可训练参数量一起报告。

这里容易混淆三种“调整模型”的方式。SFT 和 DPO 是训练目标，分别需要示范回答和成对偏好；LoRA 规定参数更新的形式，可用于不同训练目标；提示中附上范例则属于上下文学习，只改变这一次预测的条件，不改变权重。讨论某种方法的效果时，应先说清究竟改变了数据、损失还是可训练参数。

取$d=k=4,r=1,\alpha=1$，原权重有$16$个参数，$A,B$合计有八个。令$A=(1,0,0,0)$、$B=(1,2,0,0)^{\mathsf T}$、$x=(3,4,5,6)^{\mathsf T}$，则$Ax=3$，增量输出$BAx=(3,6,0,0)^{\mathsf T}$。这次更新只经由一个中间标量，能够改变的方向受限；原模型仍按$Wx$参与计算。推理时原权重与训练完成的适配器均固定。

\paragraph{低秩更新的前向与反向计算。}
对单个线性层，令输入$x\in\R^k$、输出$y\in\R^d$、冻结权重$W\in\R^{d\times k}$，并记$s=\alpha/r$。LoRA并非先构造完整更新再相乘，而可按
\begin{equation}
 v=Ax\in\R^r,\qquad y=Wx+sBv
\end{equation}
计算。若上游梯度为$g=\nabla_y\mathcal L\in\R^d$，由微分$d\mathcal L=g^{\mathsf T}dy$逐项整理得
\begin{gather}
 \nabla_B\mathcal L=s\,g(Ax)^{\mathsf T},\\
 \nabla_A\mathcal L=s\,B^{\mathsf T}g\,x^{\mathsf T},\\
 \nabla_x\mathcal L=W^{\mathsf T}g+sA^{\mathsf T}B^{\mathsf T}g.
\end{gather}
冻结$W$省去了它的参数梯度与优化器状态，却不一定省去经过该层的反向计算和激活存储。因此训练显存不会按可训练参数的比例下降。常见初始化令$A$随机、$B=0$，使起点与原模型相同。此时首步$\nabla_A=0$，$\nabla_B$通常非零；待$B$更新后，$A$也开始得到梯度。若两者都从零开始，更新无法启动。

从线性代数看，$\operatorname{rank}(BA)\leq r$，即增量最多沿$r$个独立方向改变输出。若理想的全量更新$\Delta W_*$具有奇异值$\sigma_1\geq\cdots\geq\sigma_s$，最佳秩$r$近似的 Frobenius 范数误差为$(\sum_{j>r}\sigma_j^2)^{1/2}$。这个结果说明：任务所需的变化若主要集中在少数方向，低秩约束可能足够；若重要变化分散在许多方向，小$r$便会限制性能。它是矩阵近似的结论，不保证实际训练一定找到该近似。部署时可在数值格式兼容的条件下把固定增量并入$W$，而量化或动态切换适配器还需单独验证。

\subsection{偏好优化：比较两种回答}
示范数据只告诉模型应当写出哪段回答；偏好数据则给出同一问题下的两个回答，并指出哪个更好。要理解 DPO 的目标，先考虑一个带约束的理想问题：一方面希望回答获得较高奖励，另一方面不希望新策略偏离参考模型过远。设提示为$q$，完整回答为$y$，回答概率为$\pi(y\mid q)$，参考策略为$\pi_{\rm ref}$，奖励为$r(q,y)$，相应目标为
\begin{equation}
 \max_\pi\left\{\sum_y\pi(y\mid q)r(q,y)
 -\beta\sum_y\pi(y\mid q)\log\frac{\pi(y\mid q)}{\pi_{\rm ref}(y\mid q)}\right\},\qquad\beta>0.
\end{equation}
第二项是相对于参考策略的 KL 惩罚，$\beta$控制偏离的代价。用乘子$\lambda$处理$\sum_y\pi(y\mid q)=1$的约束，对每个$\pi(y\mid q)$求导，得到
\begin{equation}
 r(q,y)-\beta\left(\log\frac{\pi(y\mid q)}{\pi_{\rm ref}(y\mid q)}+1\right)+\lambda=0.
\end{equation}
假设参考策略对被比较的回答赋予正概率，且归一化常数$Z(q)=\sum_y\pi_{\rm ref}(y\mid q)e^{r(q,y)/\beta}$有限。于是最优策略满足$\pi^*(y\mid q)=\pi_{\rm ref}(y\mid q)e^{r(q,y)/\beta}/Z(q)$。反过来，可把奖励写成策略相对参考策略的对数概率比：
\begin{equation}
 r(q,y)=\beta\log\frac{\pi^*(y\mid q)}{\pi_{\rm ref}(y\mid q)}+\beta\log Z(q).
\end{equation}
同一提示$q$下比较两个回答时，$\log Z(q)$相消。若进一步用 Bradley--Terry 模型描述偏好，即$y^+$胜过$y^-$的概率为$\sigma(r(q,y^+)-r(q,y^-))$，并用$\pi_\theta$近似上述最优策略，就得到 DPO 的训练目标：
\begin{gather}
 a_\theta=\beta\left[\log\frac{\pi_\theta(y^+\mid q)}{\pi_{\rm ref}(y^+\mid q)}-
 \log\frac{\pi_\theta(y^-\mid q)}{\pi_{\rm ref}(y^-\mid q)}\right],\\
 \mathcal L_{\rm DPO}=-\E_{(q,y^+,y^-)}\log\sigma(a_\theta),\\
 \nabla_\theta\ell_{\rm DPO}=-\beta\sigma(-a_\theta)
 [\nabla_\theta\log\pi_\theta(y^+\mid q)-\nabla_\theta\log\pi_\theta(y^-\mid q)].
\end{gather}
计算$\pi_\theta(y\mid q)$时，完整回答的对数概率是各回答词元条件对数概率之和，提示词元不计入。式中的$a_\theta$比较的是\emph{两种回答相对于参考模型的概率变化}，而非优选回答单独的绝对概率。因此 DPO 学到的是标注数据表达的偏好；若标注偏爱自信却缺乏依据的答复，模型也可能学到这种倾向。回答长度、偏好噪声与数据覆盖同样会影响结果。

例如，参考模型给两个完整回答的概率都是$0.2$，当前模型分别给出$0.4$和$0.1$。取$\beta=1$，便有$a_\theta=\log(0.4/0.2)-\log(0.1/0.2)=\log4$，偏好模型给优选回答的胜出概率为$\sigma(\log4)=0.8$，损失为$-\log0.8$。若当前模型与参考模型相同，胜出概率为$0.5$，损失为$\log2$。这里的$0.8$仅描述\emph{偏好判断}，不是该回答事实正确的概率。

\subsection{解码与推理成本}
训练完成后，模型参数固定。给定前缀，网络计算词表得分，经 softmax 得到分布，再按选定规则取出一个词元。图\ref{fig:llm-decoding-loop}中的反馈箭头表明：该词元会成为下一步前缀的一部分。训练时各位置可以并行计分，生成时却必须逐步获得尚未知晓的词元。
\begin{figure}[H]
\centering
\begin{tikzpicture}[node distance=0.72cm and 0.7cm]
 \node[idi input, minimum width=1.7cm] (ctx) {上下文\\$x_{<t}$};
 \node[idi state, right=of ctx, minimum width=1.8cm] (net) {Transformer\\隐藏状态};
 \node[idi process, right=of net, minimum width=1.8cm] (logit) {词表 logits\\$z_t$};
 \node[idi output, right=of logit, minimum width=1.8cm] (prob) {softmax\\$p_t$};
 \node[idi input, below=0.75cm of prob, minimum width=1.8cm] (sample) {采样/贪心\\选定 token};
 \draw[idi arrow] (ctx)--(net); \draw[idi arrow] (net)--(logit); \draw[idi arrow] (logit)--(prob); \draw[idi arrow] (prob)--(sample);
 \draw[idi feedback] (sample.west) -- node[below,font=\scriptsize]{追加到上下文} (ctx.south |- sample.west) -- (ctx.south);
\end{tikzpicture}
\caption{自回归解码的一步。$p_t$是给定当前前缀的词表分布；前向箭头表示计算次序，反馈箭头表示把选出的词元追加到下一步输入。}
\label{fig:llm-decoding-loop}
\end{figure}

最简单的选择规则是贪心解码，每步取概率最高的词元；束搜索则暂存若干候选序列。若希望从分布中采样，可先用温度$\tau_{\mathrm{dec}}>0$调整当前词表得分：
\begin{equation}
 p_{\tau_{\mathrm{dec}}}(x_t=j\mid x_{<t})
 =\frac{\exp(z_{t,j}/\tau_{\mathrm{dec}})}{\sum_{k=1}^{V_{\mathrm{voc}}}\exp(z_{t,k}/\tau_{\mathrm{dec}})}.
\end{equation}
这里$z_t\in\R^{V_{\mathrm{voc}}}$是当前词表的得分，$j$是候选词元编号；温度与序列长度$T$无关。温度升高，概率分布趋于平坦。采样还可以限制候选集合：top-$k$仅保留概率最高的$k$项，top-$p$则按概率降序保留累计质量首次达到$p$的最短前缀，再重新归一化。若三项概率为$0.6,0.25,0.15$，top-$2$与top-$0.8$都保留前两项，归一化后分别为$0.6/0.85$与$0.25/0.85$。$p$约束的是累计质量，而非每项的最低概率。

温度怎样改变概率，可以直接比较两个词元。固定当前得分，它们的对数概率比满足
\begin{equation}
 \log\frac{p_\tau(j)}{p_\tau(k)}=\frac{z_j-z_k}{\tau}.\label{eq:llm-temperature-ratio}
\end{equation}
由式\eqref{eq:llm-temperature-ratio}可知，$\tau\to0^+$时概率集中到最大logit集合，$\tau\to\infty$时在有限词表上趋于均匀。令$\mathcal H(\tau)=-\sum_jp_\tau(j)\log p_\tau(j)$，利用$\mathcal H=\log Z-\E_{p_\tau}[z]/\tau$及$d\E[z]/d\tau=-\Var_{p_\tau}(z)/\tau^2$，得到
\begin{equation}
 \frac{d\mathcal H}{d\tau}=\frac{\Var_{p_\tau}(z)}{\tau^3}\geq0.
\end{equation}
所以对固定前缀，升高温度会增加词表分布的熵。但生成一整段文字时，采样会改变后续前缀，不能由单步熵推出答案的事实性。top-$k$与top-$p$还会把候选集合$S$之外的概率置零，再按$p'(j)=p(j)\mathbf1_{j\in S}/\sum_{r\in S}p(r)$归一化；最终采样分布已经不同于模型原始分布。

逐步最优也未必使整段概率最大。假设第一步两条路径的概率为$0.6$与$0.4$，它们各自最好的下一步概率为$0.5$与$0.9$。完整路径概率分别是$0.30$与$0.36$，贪心解码会选择前者。束搜索保留更多路径，但束宽有限，仍是近似搜索；长度惩罚与终止规则也会影响结果。

逐词生成还涉及计算复用。KV cache 保存历史位置的键和值，避免下一步重新计算这些表示。设批大小为$B$、层数为$L$、缓存长度为$T$、键值头数为$H_{\rm kv}$、每头维度为$d_h$，每个元素占$b$字节，仅缓存键和值约需
\begin{equation}
 M_{\rm KV}=2BLTH_{\rm kv}d_hb.
\end{equation}
例如$B=1,L=2,T=1024,H_{\rm kv}=4,d_h=64,b=2$时，结果为$2097152$字节，即$2$ MiB；模型权重、当前激活和框架开销尚未计入。每新增一个词元，它的查询仍须与历史键比较，所以标准全局注意力的单步成本随已有上下文增长。缓存节省重复计算，并不让长上下文的成本保持常数。

缓存占用也随上下文增长，实际服务还要隔离不同会话并清理敏感内容。量化虽可降低内存成本，对故障码、数值、单位及工具参数的影响却应单独测试。至此讨论的解码与计算优化都只作用于已有模型；若回答需要最新规程，就要引入可更新的外部文档。

\section{检索、工具与安全边界}
\subsection{先取得适用证据，再组织回答}
预训练学到的文本分布会随语料固定；设备手册却会修订。检索增强生成（Retrieval-Augmented Generation，RAG）因此把资料查询放到回答时进行：给定问题$q$，检索器返回文档片段$D_q$，生成模型再以$q,D_q$为条件产生回答$y$\cite{lewis2020rag}。对工业规程而言，查到\emph{相关}段落还不够，必须先确认该段适用于当前型号、工况、日期与访问权限。

设人员询问“型号 M7 何时复检”。离线阶段可按章节切分手册，为每个片段保存标题、版本、页码与有效期，并建立检索索引。在线阶段先根据问题中的型号、日期及权限限定可搜索集合，再做相似度召回。双编码器分别计算问题与文档向量，文档向量可以预先缓存；交叉编码器把问题与少量候选共同输入，再精细比较。两者分别适合大范围召回与候选重排。

假设返回的一个片段只写适用条件，另一个片段虽写了处理步骤，却来自旧版手册。把两段拼接起来会产生一条原文从未给出的操作建议。此时应继续查找完整条款，找不到就说明缺少依据。图\ref{fig:llm-rag-validity}用两版阈值展示同一道理：v2 与问题在字面上相关，但已过期，只能让 v3 进入回答上下文。

图中的阈值仅为教学设定，不代表真实设备规程。实际 RAG 查询通常不会更新模型参数，只改变本次可见的文档内容。手册修订后应更新索引；编码器版本改变后，还须重算文档向量。检索片段也占上下文长度，选入太多可能挤掉真正决定适用性的条件。
\begin{figure}[H]
\centering
\begin{tikzpicture}[node distance=0.5cm, every node/.style={font=\small}]
 \node[idi input,text width=2.1cm] (query) at (0,2.2) {问题：M7 何时复检？};
 \node[idi input,text width=2.1cm] (current) at (0,0) {M7 手册 v3\\当前有效\\$80\,^{\circ}\mathrm C$ 复检};
 \node[idi input,text width=2.1cm] (old) at (0,-2.2) {M7 手册 v2\\已失效\\$75\,^{\circ}\mathrm C$ 复检};
 \node[idi process,text width=1.8cm] (filter) at (3.1,0) {型号、有效期\\和权限过滤};
 \node[idi state,text width=1.6cm] (evidence) at (6,0) {可用证据\\v3 条款};
 \node[idi output,text width=2.0cm] (answer) at (9,0) {达到 $80\,^{\circ}\mathrm C$ 复检\\引用 v3};
 \draw[idi arrow] (query.east)--(filter.north west);
 \draw[idi arrow] (current.east)--(filter.west);
 \draw[idi arrow,dashed] (old.east)--node[below,font=\scriptsize] {排除} (filter.south west);
 \draw[idi arrow] (filter.east)--(evidence.west);
 \draw[idi arrow] (evidence.east)--(answer.west);
\end{tikzpicture}
\caption{有效期约束下的手册问答示意。输入为问题与两版候选条款；实线箭头表示筛选和传递，虚线表示已失效条款被排除。阈值仅用于教学，不代表真实设备规程。}
\label{fig:llm-rag-validity}
\end{figure}

\subsection{检索得分与证据覆盖}
检索器怎样学会区分条款？设查询编码器$f_\phi$与文档编码器$g_\phi$输出归一化的$d$维向量，以内积$s_\phi(q,d)=f_\phi(q)^{\mathsf T}g_\phi(d)$打分。给定一个相关片段$d^+$与候选集合$\mathcal C$，可以用如下对比损失训练检索器：
\begin{gather}
 P_\phi(d\mid q,\mathcal C)=\frac{\exp(s_\phi(q,d)/\tau_r)}{\sum_{d'\in\mathcal C}\exp(s_\phi(q,d')/\tau_r)},\\
 \ell_{\rm ret}=-\log P_\phi(d^+\mid q,\mathcal C),\qquad \tau_r>0,\\
 \frac{\partial\ell_{\rm ret}}{\partial s_\phi(q,d)}=\frac{P_\phi(d\mid q,\mathcal C)-\mathbf1_{d=d^+}}{\tau_r}.\label{eq:llm-retrieval-gradient}
\end{gather}
式\eqref{eq:llm-retrieval-gradient}表明，若各候选得分能独立调整，下降方向会提高正例的相对得分、降低其他候选的得分。实际编码器参数由多组样本共享，更新效果须在验证集检验。更重要的是，$P_\phi(d\mid q,\mathcal C)$仅在候选集合$\mathcal C$内归一化，不是“这段文档必然支持答案”的概率。训练样本中也可能有多个正确条款，不能把另一份有效证据误标为负例。权限和有效期应先确定集合$\mathcal D(q)$，相似度只负责在其中排序。

还可以用一个\emph{理想化的概率模型}理解检索与生成的连接：假设回答依赖某一份未直接指定的文档$d$，便对所有允许文档求和，得到
\begin{equation}
 p(y\mid q)=\sum_{d\in\mathcal D(q)}p_\phi(d\mid q)p_\theta(y\mid q,d).
\end{equation}
若只能取前$K$份文档，记它们组成$S_K$，并令$Z_K=\sum_{d\in S_K}p_\phi(d\mid q)$。对留下的质量重新归一化，得到
\begin{equation}
 p_K(y\mid q)=\frac1{Z_K}\sum_{d\in S_K}p_\phi(d\mid q)p_\theta(y\mid q,d).
\end{equation}
当$0<Z_K<1$时，完整分布可写成$p(y\mid q)=Z_Kp_K(y\mid q)+(1-Z_K)p_{\rm tail}(y\mid q)$，故两者对任意回答事件的概率差不超过$1-Z_K$。这说明遗漏概率质量会影响截断近似，但只有在$p_\phi(d\mid q)$确实是全库概率分布时，$Z_K$才有上述含义。普通向量检索分数通常不满足这一条件，不能用前$K$名的 softmax 数值宣称“证据已覆盖多少”。实际把多个片段拼入同一个提示，也不同于上述单文档混合模型。

检索评价还应区分“找到证据”和“答对问题”。令$E$表示上下文包含足够的有效证据，$A$表示答案正确，则全概率公式给出
\begin{equation}
 P(A)=P(E)P(A\mid E)+P(E^{\rm c})P(A\mid E^{\rm c}).
\end{equation}
没有检索到证据，模型仍可能凭记忆答对，因此一般不能把答案正确率简单地以上述证据召回率为上界。若任务要求“正确且得到当前有效文档支持”，成功事件必然包含在$E$中，成功率才不会超过$P(E)$。这说明应分别记录证据召回、答案正确与引用支持。增大$K$可能增加召回，也可能引入冲突条款，应在固定上下文预算下比较。

在工程实现中，这些原则对应三个可检查的环节。建索引时保留设备、部件、工况、版本、页码与表格关系，不应把警告条件和处理步骤拆散。检索时可用 BM25 找精确故障码，以向量检索补充近义描述，再重排候选。生成时要求回答带上引用位置和版本；若缺少型号或工况，就先请求补充信息。

测试问题要覆盖直接查表、跨段落组合、版本冲突、过期规程、未知故障码以及恶意文本。每一类分别检查召回、答案、引用和拒答，才能辨认错误来自哪个环节。特别是版本与权限过滤应在检索之前完成，不能把所有文件都交给生成模型，再希望它自行辨别哪一份可用。

\subsection{工具调用与系统评价}
有些问题还需查询实时设备记录。此时模型生成的是\emph{工具请求}，包括工具名、设备编号和时间范围；外部程序校验格式与权限、执行查询，再把结果交回模型。请求格式正确不代表设备存在，发出请求也不代表执行成功。超时、空结果及版本冲突必须作为明确状态返回，回答不能把缺失数据补写成看似合理的数值。

工具查询发生在推理期间，不会自动改变模型参数。若工具还可能发出控制动作，应另设权限边界与人工确认。评价这类系统时，须分别检查参数是否合法、工具是否执行成功、回答是否正确使用返回值；正常请求之外，还要覆盖边界输入、过时资料和工具失败。

系统还会遇到两类性质不同的风险。其一是无证据回答：引用不存在、混淆型号，或把附条件的建议写成确定结论。其二是提示注入：检索到的文档夹带命令，企图改变模型行为。前者要求逐项核查来源并允许拒答，后者要求把检索文本视为数据而非指令；工具调用还须经过结构化参数校验、权限检查和资源范围限制。

最后将评价结果按四类汇总：语言与格式遵循、任务正确性、证据忠实度和系统可靠性。正确性应由专家标注或可验证程序判定；忠实度检查引用能否支持每项结论；可靠性记录延迟、成本、拒答、超时与工具失败。再按型号、工况、输入长度和信息缺失情况分层，才能发现平均分数掩盖的系统性错误。

\section*{本章小结}
从词元预测出发，可以看清语言模型各环节的分工。预训练拟合语料中的条件分布；示范微调与偏好优化改变回答倾向；LoRA 规定参数以何种形式更新；解码规则决定怎样从词表分布中选出文字。这些操作都没有自动赋予模型判断现行规程的能力。

工业问答因此还需要有效文档检索、明确的引用、受控的工具权限与逐项评价。涉及故障码、量值和复检条件时，回答须保留型号、版本和来源；缺少关键条件就应请求补充信息。下一章进一步讨论文字怎样与图像等观测对应：所谓“此处裂纹”究竟指哪个位置，仍须由多模态证据回答。

\section*{辅助学习材料}
以下资料可在学习相应小节后查阅；代码示例涉及的接口以实际安装版本为准。
\begingroup
\emergencystretch=3em
\begin{itemize}
\item \textbf{CS336: Language Modeling from Scratch}（斯坦福大学，英文讲义）。从分词、张量与资源核算读起，对照本章的词元表示、因果预测和 KV 缓存计算；讲义中的资源估计应结合具体模型配置使用。\href{https://github.com/stanford-cs336/lectures}{讲义仓库}。
\item \textbf{Causal language modeling}（Hugging Face，英文教程）。查看文本分块、标签构造、训练与生成示例，对照本章的下一词元目标。示例是在已有模型上继续训练。\href{https://huggingface.co/docs/transformers/tasks/language\_modeling}{在线阅读}。
\item \textbf{SFT Trainer}（Hugging Face TRL，英文教程）。重点读指令数据格式、标签右移、只对回答计损失与指令微调示例，把本章“监督微调与参数高效适配”落实为训练样本与损失掩码；再比较仅回答计分和整段文本计分的差别。\href{https://huggingface.co/docs/trl/sft\_trainer}{在线阅读}。
\item \textbf{PEFT：在低资源硬件上对十亿规模模型进行参数高效微调}（Hugging Face，中文博客）。阅读 LoRA 配置、可训练参数统计和适配器保存／加载示例，对照本章低秩更新的前向与反向计算，区分冻结基座和训练适配器。文中低精度训练示例不等于 QLoRA 的完整说明，也不意味着显存与参数按同一比例减少。\href{https://huggingface.co/blog/zh/peft}{在线阅读}。
\item \textbf{ChatGPT 背后的“功臣”——RLHF 技术详解}（Hugging Face，中文图解博客）。按预训练、奖励模型与强化学习微调三步读图，识别偏好数据和参考策略各自的作用，再与本章 DPO 推导比较：DPO 直接优化相对回答概率，并非照搬奖励模型加策略优化的训练流程。\href{https://huggingface.co/blog/zh/rlhf}{在线阅读}。
\item \textbf{一文彻底搞懂大模型 - RAG（检索、增强、生成）}（CSDN，中文博客）。阅读知识库构建、向量检索、排序与上下文融合流程，对照本章“检索增强生成”，把召回错误与生成错误分开检查；图示便于理解管线，但“增强”不保证事实正确，仍须执行本章的版本、权限和引用支持率验收。\href{https://blog.csdn.net/a2875254060/article/details/142468037}{在线阅读}。
\item \textbf{Prompt Engineering 课程}（Datawhale 收录的中文配套教程）。可用文本总结和结构化输出示例练习维修记录整理。提示仅改变当前输入，不等于微调。\href{https://github.com/datawhalechina/llm-cookbook}{在线阅读}。
\end{itemize}
\endgroup
